\chapter{Relier les outils et construire le parcours vers le master}
\section*{Exercice 30.1 -- Identifier une structure}
Expliquer le lien entre la meilleure constante pour une variable aléatoire et le point d'une droite le plus proche d'un point du plan.

\section*{Exercice 30.2 -- Deux directions propres}
Pour $A=\frac14\begin{pmatrix}3&1\\1&3\end{pmatrix}$ et $u_0=(4,0)$, trouver $u_n=A^nu_0$ et sa limite.

\section*{Exercice 30.3 -- Discret et continu}
Une suite multiplie sa valeur par $1{,}1$ à chaque unité de temps. Quel modèle $y'=ay$ reproduit exactement ces valeurs aux temps entiers ?

\section*{Exercice 30.4 -- Signal de quatre points}
Décomposer $x=(3,1,-1,1)$ sur les quatre vecteurs $v_0,v_1,v_2,v_3$ du cours.

\section*{Exercice 30.5 -- Hypothèse manquante}
Réfuter : « une fonction de dérivée positive sur son domaine est croissante » ; « des variables décorrélées sont indépendantes ».

\section*{Exercice 30.6 -- Rédiger une preuve transférable}
Prouver que, pour $n$ points $A_i$ de $\R^d$, leur moyenne $G$ minimise $\sum_i\|M-A_i\|^2$. Donner une application.
